\section{Literature Review}
\label{sec:literature}

This section reviews the relevant literature on algorithm appreciation and aversion, trust and reliance in AI, personalization in AI systems, and delegation in decision-making. We organize the literature around key themes and identify gaps that our study addresses.

\subsection{Algorithm Appreciation and Aversion}

The literature on algorithm appreciation and aversion examines when and why people prefer algorithmic judgment over human judgment, or vice versa. \citet{dietvorst2015} introduced the concept of algorithm aversion, showing that people avoid algorithms after observing them make errors, even when the algorithms outperform humans. In their study, participants were less likely to use an algorithm after seeing it make a mistake, even though the algorithm had a better track record than human forecasters.

In contrast, \citet{logg2019} found evidence of algorithm appreciation, showing that people prefer algorithmic judgment over human judgment in objective tasks. In their study, participants were more likely to trust algorithms than humans for tasks such as forecasting and numerical estimation. This suggests that algorithm aversion/appreciation depends on task type and context.

\citet{castelo2019} extended this work by showing that algorithm aversion/appreciation depends on task type (objective vs. subjective). They found that people are more likely to use algorithms for objective tasks (e.g., forecasting) than for subjective tasks (e.g., hiring). This aligns with the MABA-HABA (Machines Are Better At vs. Humans Are Better At) framework proposed by \citet{glikson2020}, which suggests that people are more willing to delegate tasks that machines are perceived to be better at.

\citet{mahmud2022} conducted a systematic review of the literature on algorithm aversion and identified several factors that influence algorithm aversion/appreciation, including task characteristics, algorithm performance, and individual differences. They found that algorithm aversion is more common in subjective tasks and when algorithms make errors, while algorithm appreciation is more common in objective tasks and when algorithms perform well. More recently, \citet{kim2026} examined the impact of psychological readiness on algorithmic advice and found that psychological readiness moderates the relationship between algorithm performance and advice utilization, while \citet{sun2026} examined whether transparency matters when an AI system meets performance expectations and found that transparency has limited effects when performance is high.

\subsection{Trust, Personalization, and Delegation}

Trust is a key factor influencing delegation to algorithms \citep{lee2004, kohn2021}. \citet{lee2004} distinguished between trust as an attitude and reliance as behavior, noting that people may trust an algorithm but still choose not to rely on it. This distinction is important for understanding delegation decisions. \citet{hoff2015} identified several factors that influence trust in automation, including perceived competence, reliability, and understandability, finding that perceived competence is the most important factor. \citet{kohn2021} emphasized the importance of measuring both trust attitudes and reliance behaviors, noting that they may not always align.

Personalization is a common feature of AI systems, from personalized recommendations on e-commerce platforms to personalized financial advice from robo-advisors. Personalization is often assumed to increase trust and adoption \citep{longoni2019, jussupow2020}. \citet{longoni2019} examined resistance to medical AI and found that personalization increases trust in medical AI, with participants more likely to trust a personalized AI that incorporated their individual characteristics than a non-personalized AI. \citet{jussupow2020} conducted a systematic review and found that personalization can reduce aversion when AI accounts for unique situations, suggesting that personalization makes AI systems more relevant and useful to individual users. \citet{ge2020} examined learning personalized risk preferences for recommendation and found that personalized algorithms can improve recommendation quality, while \citet{dean2022} examined preference dynamics under personalized recommendations and found that personalized recommendations can influence preferences over time.

However, personalization may also have unintended consequences. \citet{castelo2019} found that personalized algorithms may be perceived as less objective or less competent than algorithms that use objective criteria, and \citet{hoff2015} noted that personalization may raise concerns about privacy, surveillance, or manipulation. \citet{hu2024} examined whether too much light blinds and found that excessive transparency can reduce trust in personalized algorithms, while \citet{ning2024} examined how transparency affects algorithmic advice utilization and found that transparency has complex effects on advice utilization.

Delegation is a fundamental aspect of decision-making, both in human-human and human-AI interactions \citep{freer2023, ivanovastenzel2025}. People delegate decisions for various reasons, including lack of expertise, time constraints, or a desire to avoid responsibility \citep{freer2023}. \citet{freer2023} examined motives for delegating financial decisions and found that people delegate for reasons such as lack of expertise, time constraints, and a desire to avoid responsibility. \citet{ivanovastenzel2025} examined preferences for decision autonomy in the age of AI and found that people value decision autonomy and are reluctant to delegate decisions to algorithms.

In the context of AI, delegation decisions are influenced by factors such as perceived algorithm competence, task characteristics, and individual differences \citep{holzmeister2023, gazit2023}. \citet{holzmeister2023} examined delegation to financial advisors and found that people are more likely to delegate when they perceive the advisor as competent and trustworthy, with perceived competence being the most important factor. \citet{gazit2023} examined choosing between human and algorithmic advisors and found that people are more likely to delegate to algorithms than to human advisors for objective tasks, but not for subjective tasks. \citet{germann2022} examined algorithm aversion in delegated investing and found that people are reluctant to delegate investment decisions to algorithms, even when algorithms outperform humans, while \citet{kaawach2023} examined automatic vs. manual investing and found that past performance influences delegation decisions. \citet{renier2021} examined to err is human and found that people are more forgiving of human errors than algorithm errors, and \citet{miazek2025} examined when AI is fairer than humans and found that egocentrism influences moral and fairness judgments of AI and human decisions.

\subsection{Individual Differences and Explainability}

Individual differences such as risk attitudes, trust propensity, and technology affinity may influence delegation to algorithms \citep{dohmen2011, franke2019}. \citet{dohmen2011} examined individual risk attitudes and found that risk attitudes are stable over time and influence decision-making. \citet{franke2019} examined technology affinity and found that technology affinity influences adoption of new technologies. \citet{gaechter2022} examined individual-level loss aversion and found that loss aversion varies across individuals and influences decision-making, while \citet{koebberling2005} examined an index of loss aversion and found that loss aversion can be measured reliably. \citet{falk2016} examined a preference survey module and found that preferences can be measured reliably using survey questions, and \citet{bruin2010} examined risk and rationality and found that there is substantial heterogeneity in risk preferences.

Explainability and transparency are important factors influencing trust in AI \citep{shin2021, leightmann2023}. \citet{shin2021} examined the effects of explainability and causability on trust in human-AI collaboration and found that explainability and causability affect trust. \citet{leightmann2023} examined effects of explainable AI and found that explainability affects trust in AI. \citet{ning2024} examined how transparency affects algorithmic advice utilization and found that transparency has complex effects on advice utilization, while \citet{hu2024} examined whether too much light blinds and found that excessive transparency can reduce trust in personalized algorithms. \citet{raes2024} examined from explainable to interactive AI and found that interactive AI can improve trust and understanding, and \citet{westphal2023} examined decision control and explanations and found that decision control and explanations interact to influence trust.

\subsection{Summary and Research Gap}

The literature on algorithm appreciation and aversion, trust in AI, personalization, and delegation provides a rich foundation for understanding delegation to algorithms. However, there is a gap in our understanding of how personalization affects delegation to algorithms. While personalization is often assumed to increase trust and adoption \citep{longoni2019, jussupow2020}, this assumption has not been systematically tested in the context of delegation. Moreover, most studies on algorithm appreciation and aversion have focused on comparing algorithms to humans, rather than comparing different types of algorithms to each other \citep{dietvorst2015, logg2019}.

This study fills this gap by comparing delegation to an objective EV-maximizing algorithm versus a personalized algorithm that incorporates individual preferences. We examine how personalization affects delegation in financial decision-making, a domain where both objective criteria (expected value) and subjective preferences (risk attitudes) are relevant. Our study contributes to the literature by providing novel evidence that personalization can deter rather than encourage delegation to algorithms, challenging common assumptions about the benefits of personalization in AI systems.